\appendix

\section{Cross-validated candidate catalogue}
\label{sec:catalog}
Table~\ref{tab:catalog} lists the 49 unique DESI v3 survivors that match a Euclid Q1
discovery-engine lens, are NEW to the DESI lens-finder catalogues (Storfer/Inchausti/Huang within
$5''$), and were graded at Euclid $0.1''$ by the LensJudge v2 escalation grader. These are
independently \emph{recovered} Euclid catalogue lenses --- a four-way cross-validation (DESI v3 CNN
$\cap$ Euclid discovery engine $\cap$ Euclid expert grade $\cap$ LensJudge-at-$0.1''$) --- and are
deliberately \emph{not} claimed as net-new discoveries. The $p_{\rm lens}$ column shows the
DESI$\to$Euclid tier flip; ``release'' marks whether the object was recovered in the DR10 sweep,
the DR11-south sweep, or both.

{\small
\begin{longtable}{lcccccc}
\caption{Cross-validated strong-lens candidates: DESI v3 survivors that match a Euclid Q1 discovery-engine lens, NEW to the DESI lens-finder catalogues (Storfer/Inchausti/Huang), graded at Euclid 0.1$''$. These are independently recovered Euclid catalogue lenses, \emph{not} net-new discoveries. $p_{\rm lens}$ shows the DESI$\to$Euclid tier flip.}\\
\label{tab:catalog}\\
\toprule
DESI id & RA & Dec & Euclid & LensJudge & $p_{\rm lens}$ & release\\
 & (deg) & (deg) & grade & @0.1$''$ & DESI$\to$Euclid & \\
\midrule
\endfirsthead
\toprule DESI id & RA & Dec & Euclid & LensJudge & $p_{\rm lens}$ & release\\\midrule\endhead
s\_83157\_8409 & 57.7025 & -48.4062 & A & A & 0.12$\to$0.97 & both\\
s\_85094\_13993 & 63.6294 & -48.0333 & A & A & 0.30$\to$0.97 & DR10\\
s\_79373\_7769 & 58.2019 & -49.4705 & A & A & 0.05$\to$0.95 & both\\
s\_80329\_6402 & 64.5419 & -49.3086 & A & A & 0.03$\to$0.95 & both\\
s\_81272\_7726 & 63.9239 & -48.9296 & A & A & 0.03$\to$0.93 & both\\
s\_86048\_6517 & 58.6493 & -47.6453 & A & A & 0.02$\to$0.93 & DR10\\
s\_169188\_5518 & 53.3251 & -29.1510 & A & A & 0.15$\to$0.93 & both\\
s\_179335\_8072 & 51.0363 & -27.2415 & A & A & 0.01$\to$0.93 & both\\
s\_84136\_4630 & 65.1956 & -48.3600 & A & A & 0.07$\to$0.88 & DR10\\
s\_174235\_13613 & 51.5787 & -28.1897 & A & A & 0.02$\to$0.88 & both\\
s\_178058\_11508 & 51.9633 & -27.4260 & A & A & 0.04$\to$0.88 & both\\
s\_165427\_7416 & 52.9210 & -29.9100 & A & A & 0.03$\to$0.87 & both\\
s\_175513\_5175 & 53.5621 & -28.0129 & A & A & 0.02$\to$0.87 & DR11\\
s\_82216\_7491 & 62.1645 & -48.8224 & A & A & 0.09$\to$0.85 & both\\
s\_93968\_433 & 62.8684 & -45.7047 & A & A & 0.04$\to$0.84 & both\\
s\_170450\_15146 & 54.4805 & -28.9549 & A & A & 0.04$\to$0.82 & both\\
s\_81263\_7090 & 60.4836 & -49.0669 & A & B & 0.15$\to$0.73 & both\\
s\_171714\_3810 & 54.7757 & -28.6645 & A & B & 0.02$\to$0.72 & DR11\\
s\_90958\_9241 & 59.7052 & -46.5695 & A & B & 0.03$\to$0.70 & both\\
s\_176778\_4403 & 50.8798 & -27.6486 & A & B & 0.03$\to$0.70 & both\\
s\_170452\_14504 & 54.9726 & -29.0757 & A & B & 0.18$\to$0.65 & both\\
s\_84129\_16461 & 62.7892 & -48.2719 & A & B & 0.04$\to$0.62 & both\\
s\_88015\_3655 & 65.1247 & -47.3087 & A & B & 0.05$\to$0.62 & both\\
s\_88988\_8317 & 62.3367 & -47.0669 & A & B & 0.03$\to$0.62 & both\\
s\_167928\_8112 & 52.3902 & -29.5927 & B & A & 0.02$\to$0.95 & both\\
s\_84129\_17659 & 62.8043 & -48.2286 & B & A & 0.03$\to$0.93 & DR10\\
s\_85093\_15882 & 63.2503 & -48.0737 & B & A & 0.04$\to$0.93 & both\\
s\_77522\_7961 & 63.7748 & -50.0119 & B & A & 0.02$\to$0.91 & both\\
s\_180619\_719 & 51.2423 & -27.0221 & B & A & 0.02$\to$0.83 & both\\
s\_179345\_702 & 53.6441 & -27.2233 & B & B & 0.03$\to$0.74 & both\\
s\_86061\_15501 & 63.6862 & -47.8550 & B & B & 0.03$\to$0.72 & both\\
s\_166673\_9358 & 51.7071 & -29.7125 & B & B & 0.03$\to$0.72 & both\\
s\_77510\_6177 & 59.0844 & -49.9562 & B & B & 0.02$\to$0.65 & DR10\\
s\_80314\_12621 & 59.0815 & -49.3650 & B & B & 0.04$\to$0.65 & both\\
s\_78453\_11482 & 63.9706 & -49.7278 & B & B & 0.02$\to$0.65 & both\\
s\_84123\_2953 & 60.2940 & -48.3238 & B & B & 0.03$\to$0.65 & both\\
s\_88997\_602 & 65.3650 & -47.0251 & B & B & 0.04$\to$0.65 & both\\
s\_81263\_7679 & 60.4929 & -48.9067 & B & B & 0.03$\to$0.62 & DR10\\
s\_87026\_5695 & 60.5672 & -47.6155 & B & B & 0.02$\to$0.62 & DR10\\
s\_88007\_7987 & 62.2189 & -47.2949 & B & B & 0.03$\to$0.62 & DR10\\
s\_84125\_13400 & 61.3553 & -48.2723 & B & B & 0.04$\to$0.60 & DR11\\
s\_88996\_2557 & 65.0706 & -47.0601 & B & B & 0.02$\to$0.60 & DR11\\
s\_167935\_8040 & 54.3727 & -29.3848 & B & B & 0.04$\to$0.60 & DR11\\
s\_75664\_6939 & 58.9597 & -50.5340 & B & B & 0.04$\to$0.55 & DR11\\
s\_81254\_8721 & 57.2304 & -49.0460 & B & B & 0.03$\to$0.55 & DR11\\
s\_88990\_5749 & 63.0160 & -47.0193 & B & B & 0.03$\to$0.55 & both\\
s\_176778\_10456 & 51.0181 & -27.7804 & B & B & 0.03$\to$0.55 & both\\
s\_76594\_8452 & 62.6149 & -50.3258 & B & B & 0.04$\to$0.50 & DR11\\
s\_176791\_230 & 54.4517 & -27.8613 & C & A & 0.03$\to$0.93 & both\\
\bottomrule
\end{longtable}
}


\section{Panel-qualified NEW grade-A candidates}
\label{sec:qualified}
Table~\ref{tab:qualified} lists the NEW grade-A candidates surfaced by v3 in the DR10 vetting and
their verdict under the campaign qualification panel (advocate/skeptic/morphology/contaminant, with
a skeptic veto and a 2-of-$N$ A-rule --- the same mechanism that took the v2 campaign's first-pass
A/B from 8 to 5). All four survive at $\ge$B with no contaminant flag, which is cleaner than the v2
campaign. These positions have \emph{no} Euclid coverage, however, so they are DESI-resolution-
qualified candidates, not confirmations.

\begin{table}[t]\centering
\caption{v3 grade-A C-vet survivors that are \emph{absent} from the Huang/Storfer/Inchausti lineage catalogues, with their literature-wide NED/SIMBAD crossmatch (\texttt{94\_external\_xmatch.py}). All three are previously-published lenses from other surveys --- v3 \emph{rediscoveries}, not net-new. (A fourth C-vet survivor, \texttt{s\_310364\_6649}, is the Huang+2020 lens \texttt{DESI-038.2078-03.3906}.) Net-new lenses: zero; the result is independent external validation of the v3 finder.}
\label{tab:qualified}\small
\begin{tabular}{lcclc}
\toprule
DESI id & C-vet $p_{\rm lens}$ & ranking & known lens (survey) & sep\\
\midrule
s\_441355\_1111 & 0.74 & head & CSWA 1 (CASSOWARY) & 0.5$''$\\
s\_124958\_10481 & 0.82 & v3blend8 & DES J035242.40-382544.9 (DES) & 0.2$''$\\
s\_102952\_9916 & 0.76 & v3blend8 & DES J221912.39-434835.1 (DES) & 0.2$''$\\
\bottomrule
\end{tabular}
\end{table}


\section{Reproducibility}
\label{sec:repro}
All results are reproduced from tracked code and artifacts on branch
\texttt{reproductions/claudenet-v3}. The script numbering is: \texttt{00--91} v1 (members, calibration,
combiners, the seven-phase program, tables/figures); \texttt{100--165} v2 (NegEval-1M, mining,
refit, the DR9 sweep, group-conformal selection); \texttt{200--260} the DR9 qualification campaign;
and \texttt{300--367} v3 --- \texttt{300--326} workstream A (mimic metric, bank, retraining,
ensemble refit, lens-vs-mimic head), \texttt{360--367} deployment (DR10/DR11 parents, scoring,
selection, the Euclid cross-validation and the model-progression comparison). The vetting harness
lives under \texttt{reproductions/lensjudge}. Shared libraries (\texttt{\_ensemble},
\texttt{\_train}, \texttt{\_clib}, \ldots) reuse the baseline matched-FPR arithmetic verbatim, so
every v1/v2/v3 number is directly comparable. The v3 figures and tables in this report are
regenerated by \texttt{92\_make\_v3\_figures.py} and \texttt{92\_make\_v3\_tables.py} from the
artifacts under \texttt{data/v3/} and \texttt{v3/}.

Survey-scale compute ran on NERSC Perlmutter (GPU account \texttt{cosmo\_g}, CPU \texttt{cosmo}); the
v3 program used $\mathcal{O}(10^2)$ GPU-hours. Agentic vetting (LensJudge) spent $\sim$\$53 of a \$100
cap, targeted on the most-promising candidates with an evidence gate. We carry the lineage's honesty
discipline throughout: ``candidate'' never ``discovery''; recall of the published catalogues is
reported before any new-candidate count; both independent graders are Claude models (different
harness/prompt, not statistically independent); the mimic-FPR seed metric is selection-biased and
its cross-model reading is corrected against the independent Euclid measure (\S\ref{sec:v3limits});
and DR11 results are reported embargo-aware (internal until the DR11 public release).
